Sources : Insee, indice Notaires-Insee des prix des logements anciens (France métropolitaine, série brute) ; indice du coût de la construction (ICC) ; indice des prix à la consommation, base 2015, loyers effectivement payés (04.1.1) et ensemble hors tabac.\par Lecture : de 1998 à 2025, le prix des logements anciens a été multiplié par 3,0, le coût de la construction par 2,0 et les loyers payés par 1,5. L'indice des loyers couvre tout le parc loué, logement social et baux en cours compris : il progresse moins vite que les loyers de relocation dans le secteur libre.
